% =============================================================================
% Model Architecture Section — HHRoad
% Hyperbolic Hierarchical Road Network Representation Learning
% =============================================================================

\section{Methodology}\label{sec:method}

In this section, we present HHRoad, a hyperbolic hierarchical road network representation learning framework.
We first introduce the Lorentz hyperbolic space and its key operations (\S\ref{subsec:prelim}),
then describe the complete model pipeline:
from raw road features to Euclidean embeddings (\S\ref{subsec:feat_embed}),
from Euclidean space to the Lorentz manifold (\S\ref{subsec:hyp_proj}),
through tangent-space graph convolution (\S\ref{subsec:hyp_gcn}),
hierarchical message passing (\S\ref{subsec:hier_mp}),
and finally the learning objectives that jointly enforce entailment constraints and contrastive alignment (\S\ref{subsec:loss}).
Figure~\ref{fig:architecture} illustrates the overall architecture.

% -------------------------------------------------------
\subsection{Lorentz Hyperbolic Space}\label{subsec:prelim}
% -------------------------------------------------------

Road networks exhibit natural hierarchical structures:
individual road \emph{segments} compose \emph{localities} (functional clusters),
which in turn form higher-level \emph{regions} (administrative zones).
Euclidean space distorts such tree-like hierarchies because the volume of a Euclidean ball grows only polynomially with its radius,
whereas hyperbolic space grows \emph{exponentially}, thereby providing ample capacity to embed hierarchies with low distortion~\cite{nickel2017poincare,ganea2018hyperbolic}.

We adopt the \textbf{Lorentz model} (also known as the hyperboloid model) of $d$-dimensional hyperbolic space.
The Lorentz manifold $\mathbb{H}^d$ is defined as
\begin{equation}\label{eq:lorentz_def}
  \mathbb{H}^{d} = \bigl\{\, \mathbf{x} = (x_0, x_1, \dots, x_d) \in \mathbb{R}^{d+1}
  \;\big|\; \langle \mathbf{x}, \mathbf{x} \rangle_{\mathcal{L}} = -1,\; x_0 > 0 \,\bigr\},
\end{equation}
where the \textbf{Minkowski inner product} is
\begin{equation}\label{eq:minkowski}
  \langle \mathbf{x}, \mathbf{y} \rangle_{\mathcal{L}}
  = -x_0 y_0 + \sum_{i=1}^{d} x_i y_i.
\end{equation}
The \textbf{geodesic distance} between two points $\mathbf{x}, \mathbf{y} \in \mathbb{H}^d$ is
\begin{equation}\label{eq:lorentz_dist}
  d_{\mathbb{H}}(\mathbf{x}, \mathbf{y})
  = \operatorname{arcosh}\!\bigl(-\langle \mathbf{x}, \mathbf{y} \rangle_{\mathcal{L}}\bigr).
\end{equation}
The \textbf{origin} of the Lorentz model is $\mathbf{o} = (1, 0, \dots, 0) \in \mathbb{H}^d$.
At any point $\mathbf{x} \in \mathbb{H}^d$, the \textbf{tangent space} $\mathcal{T}_{\mathbf{x}}\mathbb{H}^d$ consists of all vectors $\mathbf{v} \in \mathbb{R}^{d+1}$ satisfying $\langle \mathbf{x}, \mathbf{v} \rangle_{\mathcal{L}} = 0$.
Moving between the tangent space and the manifold is accomplished via the \textbf{exponential map} and the \textbf{logarithmic map}:
\begin{align}
  \exp_{\mathbf{x}}(\mathbf{v})
  &= \cosh\!\bigl(\|\mathbf{v}\|_{\mathcal{L}}\bigr)\,\mathbf{x}
     + \sinh\!\bigl(\|\mathbf{v}\|_{\mathcal{L}}\bigr)\,
       \frac{\mathbf{v}}{\|\mathbf{v}\|_{\mathcal{L}}}, \label{eq:exp_map} \\[4pt]
  \log_{\mathbf{x}}(\mathbf{y})
  &= \frac{d_{\mathbb{H}}(\mathbf{x},\mathbf{y})}
          {\sinh\!\bigl(d_{\mathbb{H}}(\mathbf{x},\mathbf{y})\bigr)}\,
     \bigl(\mathbf{y} + \langle \mathbf{x}, \mathbf{y} \rangle_{\mathcal{L}}\,\mathbf{x}\bigr),
  \label{eq:log_map}
\end{align}
where $\|\mathbf{v}\|_{\mathcal{L}} = \sqrt{\langle \mathbf{v}, \mathbf{v} \rangle_{\mathcal{L}}}$ denotes the Lorentzian norm of a tangent vector.
These two maps are central to every subsequent component, as they allow us to perform efficient linear operations in the tangent space while preserving the geometric properties of hyperbolic space.

% -------------------------------------------------------
\subsection{Road Feature Embedding}\label{subsec:feat_embed}
% -------------------------------------------------------

We consider a road network $\mathcal{G} = (\mathcal{V}, \mathcal{E})$
with $N = |\mathcal{V}|$ road segments and adjacency matrix $\mathbf{A} \in \{0,1\}^{N \times N}$.
Each segment $v_i$ is described by four categorical attributes:
node identity, road type, discretised length, and lane count.
We embed each attribute via a learnable embedding table and concatenate the results:
\begin{equation}\label{eq:raw_feat}
  \mathbf{e}_i
  = \bigl[\,
      \mathbf{E}_{\text{node}}(v_i) \;\|\;
      \mathbf{E}_{\text{type}}(v_i) \;\|\;
      \mathbf{E}_{\text{len}}(v_i)  \;\|\;
      \mathbf{E}_{\text{lane}}(v_i)
    \,\bigr]
  \;\in\; \mathbb{R}^{D_e},
\end{equation}
where $\|$ denotes concatenation and $D_e = d_{\text{node}} + d_{\text{type}} + d_{\text{len}} + d_{\text{lane}}$ is the total Euclidean feature dimension.
We further assume two \emph{assignment matrices} are given: a \textbf{structural assignment} $\mathbf{S} \in \mathbb{R}^{N \times M}$
mapping segments to $M$ localities (obtained via spectral clustering on the topology), and a \textbf{functional assignment} $\mathbf{F} \in \mathbb{R}^{M \times R}$
mapping localities to $R$ regions (obtained via trajectory-data augmented clustering).
These assignments define the three-level hierarchy:
\emph{Segment} $\rightarrow$ \emph{Locality} $\rightarrow$ \emph{Region}.

% -------------------------------------------------------
\subsection{Projection to the Lorentz Manifold}\label{subsec:hyp_proj}
% -------------------------------------------------------

To operate in hyperbolic space, we project the Euclidean features $\{\mathbf{e}_i\}_{i=1}^N$ onto the Lorentz manifold $\mathbb{H}^d$ via a learnable linear transformation followed by Lorentz projection.
Specifically, for each segment $v_i$:
\begin{align}
  \mathbf{z}_i &= \mathbf{W}_p\,\mathbf{e}_i + \mathbf{b}_p
    \;\in\; \mathbb{R}^{d}, \label{eq:linear_proj}\\[4pt]
  \mathbf{h}_i^{(0)} &= \pi_{\mathbb{H}}(\mathbf{z}_i)
    = \Bigl(\sqrt{1 + \|\mathbf{z}_i\|^2},\; \mathbf{z}_i\Bigr)
    \;\in\; \mathbb{H}^{d}, \label{eq:lorentz_proj}
\end{align}
where $\mathbf{W}_p \in \mathbb{R}^{d \times D_e}$, $\mathbf{b}_p \in \mathbb{R}^{d}$ are learnable parameters, and $\pi_{\mathbb{H}}$ prepends the time component $x_0 = \sqrt{1 + \|\mathbf{z}_i\|^2}$ so that the resulting $(d{+}1)$-dimensional vector satisfies the Lorentz constraint $\langle \mathbf{h}_i^{(0)}, \mathbf{h}_i^{(0)} \rangle_{\mathcal{L}} = -1$.
This gives us the initial hyperbolic representations $\mathbf{H}^{(0)} = \{\mathbf{h}_1^{(0)}, \dots, \mathbf{h}_N^{(0)}\} \subset \mathbb{H}^d$ for all segments.

% -------------------------------------------------------
\subsection{Hyperbolic Graph Convolution}\label{subsec:hyp_gcn}
% -------------------------------------------------------

Standard graph convolution cannot be directly applied in hyperbolic space because the manifold is non-linear.
We adopt a \textbf{tangent-space strategy}~\cite{chami2019hyperbolic}:
map features to the tangent space at the origin $\mathbf{o}$, perform linear neighbourhood aggregation there, and then map the result back to the manifold.
Concretely, given hyperbolic node features $\mathbf{H} = \{\mathbf{h}_i\}_{i=1}^N \subset \mathbb{H}^d$ and adjacency $\mathbf{A}$, a single \textbf{Hyperbolic Graph Convolution} (HGC) layer operates as follows:

\paragraph{Step 1: Map to tangent space.}
For each node $i$, extract the spatial component $\mathbf{h}_i^s = (h_{i,1}, \dots, h_{i,d})$ and compute the tangent-space representation via the logarithmic map at the origin:
\begin{equation}\label{eq:log_zero}
  \tilde{\mathbf{h}}_i
  = \log_{\mathbf{o}}^{0}(\mathbf{h}_i^s)
  = \frac{\operatorname{artanh}(\|\mathbf{h}_i^s\|)}{\|\mathbf{h}_i^s\|}\,\mathbf{h}_i^s
  \;\in\; \mathbb{R}^{d}.
\end{equation}

\paragraph{Step 2: Neighbourhood aggregation.}
Aggregate neighbour features using degree-normalised message passing:
\begin{equation}\label{eq:agg}
  \bar{\mathbf{h}}_i
  = \sum_{j \in \mathcal{N}(i)} \frac{A_{ij}}{\deg(i)}\,\tilde{\mathbf{h}}_j
  \;\in\; \mathbb{R}^{d},
\end{equation}
where $\mathcal{N}(i) = \{j : A_{ij} > 0\}$ and $\deg(i) = \sum_j A_{ij}$.
This aggregation is an Euclidean operation in the tangent space, hence efficient and numerically stable.

\paragraph{Step 3: Linear transformation.}
Apply a learnable weight matrix and bias:
\begin{equation}\label{eq:tangent_linear}
  \hat{\mathbf{h}}_i = \mathbf{W}\,\bar{\mathbf{h}}_i + \mathbf{b}
  \;\in\; \mathbb{R}^{d},
\end{equation}
with $\mathbf{W} \in \mathbb{R}^{d \times d}$ and $\mathbf{b} \in \mathbb{R}^d$.

\paragraph{Step 4: Map back to the manifold.}
Apply the exponential map at the origin to return the transformed features to hyperbolic space, and re-project to guarantee the Lorentz constraint:
\begin{align}
  \hat{\mathbf{h}}_i^{s}
  &= \exp_{\mathbf{o}}^{0}(\hat{\mathbf{h}}_i)
   = \frac{\tanh(\|\hat{\mathbf{h}}_i\|)}{\|\hat{\mathbf{h}}_i\|}\,\hat{\mathbf{h}}_i
  \;\in\; \mathbb{R}^{d}, \label{eq:exp_zero}\\[4pt]
  \mathbf{h}_i'
  &= \pi_{\mathbb{H}}\!\bigl(\hat{\mathbf{h}}_i^{s}\bigr)
  \;\in\; \mathbb{H}^{d}. \label{eq:reproj}
\end{align}
We denote one full pass of Eqs.~\eqref{eq:log_zero}--\eqref{eq:reproj} as $\operatorname{HGC}(\mathbf{H}, \mathbf{A})$.

% -------------------------------------------------------
\subsection{Hierarchical Hyperbolic Message Passing}\label{subsec:hier_mp}
% -------------------------------------------------------

A single HGC layer captures only local neighbourhood information.
To propagate information across the three-level hierarchy, we design a \textbf{Hierarchical Hyperbolic Message Passing} (HHMP) module
that first aggregates representations \emph{bottom-up} from segments to localities and regions,
and then distributes global context \emph{top-down} back to segments.
Two HHMP layers are stacked to produce the final representations.

\subsubsection{Bottom-Up Aggregation}\label{subsubsec:bottom_up}

Given segment-level hyperbolic embeddings $\mathbf{H}^{\text{seg}} = \{\mathbf{h}_i\}_{i=1}^N \subset \mathbb{H}^d$, we compute locality-level and region-level embeddings
by performing weighted Fréchet means in the tangent space at the origin.

\paragraph{Segment $\to$ Locality.}
Let $\hat{\mathbf{S}} \in \mathbb{R}^{N \times M}$ be the column-normalised structural assignment matrix.
For each locality $\ell \in \{1, \dots, M\}$:
\begin{align}
  \tilde{\mathbf{h}}_i &= \log_{\mathbf{o}}(\mathbf{h}_i),\quad i = 1, \dots, N, \label{eq:seg_to_tangent}\\[3pt]
  \bar{\mathbf{t}}_\ell &= \sum_{i=1}^{N} \hat{S}_{i\ell}\,\tilde{\mathbf{h}}_i, \label{eq:tangent_mean_loc}\\[3pt]
  \mathbf{h}_\ell^{\text{loc}} &= \pi_{\mathbb{H}}\!\Bigl(
    \rho \cdot \frac{\|\bar{\mathbf{t}}_\ell^s\|}{\|\exp_{\mathbf{o}}(\bar{\mathbf{t}}_\ell)^s\|}
    \cdot \exp_{\mathbf{o}}(\bar{\mathbf{t}}_\ell)^s
  \Bigr), \label{eq:loc_emb}
\end{align}
where the superscript $s$ extracts the spatial component and $\rho \in (0,1]$ is a \textbf{norm-preservation factor} (set to $0.8$) that re-scales the spatial norm after aggregation.
This rescaling is critical:
because the Fréchet mean in negatively-curved space contracts towards the origin, directly averaging in the tangent space would cause the locality embedding to collapse near $\mathbf{o}$.
By retaining a fraction $\rho$ of the expected spatial norm, we preserve representational capacity while still allowing a controlled norm reduction that reflects the coarser granularity of localities.

\paragraph{Locality $\to$ Region.}
The same procedure is applied to the locality embeddings $\mathbf{H}^{\text{loc}}$ using the functional assignment $\hat{\mathbf{F}}$, yielding region embeddings $\mathbf{H}^{\text{reg}} = \{\mathbf{h}_r^{\text{reg}}\}_{r=1}^R \subset \mathbb{H}^d$.

\subsubsection{Top-Down Message Passing}\label{subsubsec:top_down}

After obtaining multi-scale representations, context is propagated downward through five sequential stages,
mirroring the nomenclature of the original HRNR~\cite{hrnr}:

\paragraph{(i) Region-to-Region (F2F).}
We first enrich region representations by exchanging messages among regions.
A hyperbolic affinity matrix is constructed by
\begin{equation}\label{eq:region_affinity}
  \alpha_{r r'} = \exp\!\bigl(-d_{\mathbb{H}}(\mathbf{h}_r^{\text{reg}}, \mathbf{h}_{r'}^{\text{reg}})\bigr),
\end{equation}
and the region embeddings are updated via an HGC layer:
\begin{equation}\label{eq:f2f}
  \mathbf{H}^{\text{reg}} \leftarrow \operatorname{HGC}_{\text{reg}}\!\bigl(\mathbf{H}^{\text{reg}},\, \boldsymbol{\alpha} + \mathbf{I}\bigr).
\end{equation}

\paragraph{(ii) Region-to-Locality (F2C).}
Region-level context is distributed to each locality by the inverse of the tangent-space aggregation.
For locality $\ell$, the region message is computed as
\begin{equation}\label{eq:f2c}
  \mathbf{m}_\ell^{\text{reg}}
  = \pi_{\mathbb{H}}\!\Biggl(
      \rho' \cdot \exp_{\mathbf{o}}\!\Bigl(
        \sum_{r=1}^{R} \frac{F_{\ell r}}{\sum_{r'} F_{\ell r'}}\;
        \log_{\mathbf{o}}\!\bigl(\mathbf{h}_r^{\text{reg}}\bigr)
      \Bigr)^s
    \Biggr),
\end{equation}
where $\rho' = 0.9$ is a higher norm-preservation rate for top-down distribution (since coarse-grained information should be preserved more faithfully when mapped to finer granularity).
The locality embedding is then updated via a \textbf{gated hyperbolic interpolation}:
\begin{equation}\label{eq:gate_loc}
  \mathbf{h}_\ell^{\text{loc}}
  \leftarrow \exp_{\mathbf{h}_\ell^{\text{loc}}}\!\Bigl(
    g_\ell \cdot \log_{\mathbf{h}_\ell^{\text{loc}}}\!\bigl(\mathbf{m}_\ell^{\text{reg}}\bigr)
  \Bigr),
\end{equation}
where $g_\ell = \sigma\!\bigl(\mathbf{w}_c^\top [\mathbf{h}_\ell^{\text{loc}} \| \mathbf{m}_\ell^{\text{reg}}] + b_c\bigr) \in (0,1)$ is a learned gate
that controls how much region-level information is incorporated.
Intuitively, the logarithmic map $\log_{\mathbf{h}_\ell^{\text{loc}}}(\mathbf{m}_\ell^{\text{reg}})$ computes a tangent direction from the current locality embedding towards the region message,
the gate $g_\ell$ scales this direction,
and the exponential map moves along the geodesic by the scaled amount.

\paragraph{(iii) Locality-to-Locality (C2C).}
Updated locality embeddings exchange messages within a pre-computed structural adjacency $\mathbf{A}^{\text{loc}}$:
\begin{equation}\label{eq:c2c}
  \mathbf{H}^{\text{loc}} \leftarrow \operatorname{HGC}_{\text{loc}}\!\bigl(\mathbf{H}^{\text{loc}},\, \mathbf{A}^{\text{loc}}\bigr).
\end{equation}

\paragraph{(iv) Locality-to-Segment (C2N).}
Locality context is distributed to individual segments analogously to Eq.~\eqref{eq:f2c}, using the structural assignment $\mathbf{S}$ and a segment-level gate:
\begin{equation}\label{eq:c2n_msg}
  \mathbf{m}_i^{\text{loc}}
  = \pi_{\mathbb{H}}\!\Biggl(
      \rho' \cdot \exp_{\mathbf{o}}\!\Bigl(
        \sum_{\ell=1}^{M} \frac{S_{i\ell}}{\sum_{\ell'} S_{i\ell'}}\;
        \log_{\mathbf{o}}\!\bigl(\mathbf{h}_\ell^{\text{loc}}\bigr)
      \Bigr)^s
    \Biggr),
\end{equation}
\begin{equation}\label{eq:gate_seg}
  \mathbf{h}_i
  \leftarrow \exp_{\mathbf{h}_i}\!\Bigl(
    g_i \cdot \log_{\mathbf{h}_i}\!\bigl(\mathbf{m}_i^{\text{loc}}\bigr)
  \Bigr),
\end{equation}
with $g_i = \sigma\!\bigl(\mathbf{w}_s^\top [\mathbf{h}_i \| \mathbf{m}_i^{\text{loc}}] + b_s\bigr)$.

\paragraph{(v) Segment-to-Segment (N2N).}
Finally, segment embeddings are refined by local topology:
\begin{equation}\label{eq:n2n}
  \mathbf{H}^{\text{seg}} \leftarrow \operatorname{HGC}_{\text{seg}}\!\bigl(\mathbf{H}^{\text{seg}},\, \mathbf{A}\bigr).
\end{equation}

After passing through two stacked HHMP layers, we denote the resulting segment embeddings as
$\mathbf{H}^{\text{seg}} = \{\mathbf{h}_i\}_{i=1}^{N}$,
locality embeddings as $\mathbf{H}^{\text{loc}} = \{\mathbf{h}_\ell^{\text{loc}}\}_{\ell=1}^{M}$,
and region embeddings as $\mathbf{H}^{\text{reg}} = \{\mathbf{h}_r^{\text{reg}}\}_{r=1}^{R}$.

\subsubsection{Output Representation}\label{subsubsec:output}

The final road-segment representation is obtained by concatenating the initial hyperbolic features $\mathbf{h}_i^{(0)}$ with the refined segment embedding $\mathbf{h}_i$,
and projecting back to a $D_o$-dimensional Euclidean vector via a linear layer:
\begin{equation}\label{eq:output}
  \mathbf{r}_i = \mathbf{W}_o\,[\mathbf{h}_i^{(0)} \| \mathbf{h}_i] + \mathbf{b}_o
  \;\in\; \mathbb{R}^{D_o},
\end{equation}
where $\mathbf{W}_o \in \mathbb{R}^{D_o \times 2(d+1)}$ and $\mathbf{b}_o \in \mathbb{R}^{D_o}$.
This Euclidean output can be fed to any downstream task head.

% -------------------------------------------------------
\subsection{Learning Objectives}\label{subsec:loss}
% -------------------------------------------------------

Training is driven by three complementary objectives.

\subsubsection{Road Function Classification Loss}\label{subsubsec:cls_loss}

Following HRNR, we use a road-function classification task as the primary supervision signal.
Given the output representation $\mathbf{r}_i$ from Eq.~\eqref{eq:output},
the classification loss is the standard cross-entropy:
\begin{equation}\label{eq:loss_cls}
  \mathcal{L}_{\text{cls}}
  = -\frac{1}{|\mathcal{B}|} \sum_{i \in \mathcal{B}}
    \log \frac{\exp(\mathbf{r}_i^\top \mathbf{w}_{y_i})}
         {\sum_{c=1}^{C} \exp(\mathbf{r}_i^\top \mathbf{w}_{c})},
\end{equation}
where $\mathcal{B}$ is the training batch and $y_i$ is the ground-truth label of segment $v_i$.

\subsubsection{Entailment Cone Loss}\label{subsubsec:entail_loss}

To explicitly encode hierarchical containment in hyperbolic space,
we introduce \textbf{entailment cones}~\cite{ganea2018hyperbolic,pal2024hycoclip}.
At any point $\mathbf{x} \in \mathbb{H}^d$, an entailment cone is defined by a half-aperture angle
\begin{equation}\label{eq:aperture}
  \theta(\mathbf{x})
  = 2\arcsin\!\Bigl(\frac{1}{\cosh\bigl(d_{\mathbb{H}}(\mathbf{x}, \mathbf{o})\bigr)}\Bigr).
\end{equation}
Points closer to the origin $\mathbf{o}$ have a wider aperture and thus can ``entail'' (contain) more child concepts;
points farther from the origin are more specific and have a narrower aperture.
Given a parent point $\mathbf{x}$ and a candidate child point $\mathbf{y}$,
we compute the angle between them as
\begin{equation}\label{eq:angle}
  \phi(\mathbf{x}, \mathbf{y})
  = \arccos\!\Biggl(
    \frac{\langle \mathbf{x}, \mathbf{y} \rangle_{\mathcal{L}}}
         {\|\mathbf{x}\|_{\mathcal{L}} \cdot \|\mathbf{y}\|_{\mathcal{L}}}
    \Biggr),
\end{equation}
and define the \textbf{entailment score}:
\begin{equation}\label{eq:entail_score}
  s(\mathbf{x}, \mathbf{y}) = \theta(\mathbf{x}) - \phi(\mathbf{x}, \mathbf{y}).
\end{equation}
A positive score indicates that $\mathbf{y}$ falls inside the entailment cone of $\mathbf{x}$, i.e., $\mathbf{x}$ entails $\mathbf{y}$.

We enforce three types of entailment relations mirroring the road-network hierarchy:
\begin{enumerate}[leftmargin=*,itemsep=2pt]
  \item \textbf{Region $\Rightarrow$ Locality}:
    for every region $r$ and locality $\ell$ assigned to $r$;
  \item \textbf{Locality $\Rightarrow$ Segment}:
    for every locality $\ell$ and segment $i$ assigned to $\ell$;
  \item \textbf{Segment $\Leftrightarrow$ Segment}:
    for every pair of topologically adjacent segments $(i,j) \in \mathcal{E}$.
\end{enumerate}
The entailment loss penalises violations via a hinge function:
\begin{equation}\label{eq:loss_ce}
  \mathcal{L}_{\text{ent}}
  = \frac{1}{|\mathcal{P}|}
    \sum_{(\mathbf{x},\mathbf{y}) \in \mathcal{P}}
    \operatorname{ReLU}\!\bigl(-s(\mathbf{x}, \mathbf{y})\bigr),
\end{equation}
where $\mathcal{P}$ is the set of all parent--child pairs sampled from the three types above.

\subsubsection{Hierarchical Contrastive Loss}\label{subsubsec:contrast_loss}

To further sharpen the embeddings, we apply contrastive learning at two granularities.

\paragraph{Segment-level contrastive learning.}
For each sampled edge $(i,j) \in \mathcal{E}$, the anchor is segment $i$ and its positive is the adjacent segment $j$.
We draw $K$ negative segments uniformly at random.
The similarity is defined as the negated Lorentz distance, and the loss follows the InfoNCE formulation:
\begin{equation}\label{eq:loss_seg_cl}
  \mathcal{L}_{\text{seg}}
  = -\frac{1}{|\mathcal{E}_s|}
    \sum_{(i,j) \in \mathcal{E}_s}
    \log \frac{
      \exp\!\bigl(-d_{\mathbb{H}}(\mathbf{h}_i, \mathbf{h}_j) / \tau\bigr)
    }{
      \exp\!\bigl(-d_{\mathbb{H}}(\mathbf{h}_i, \mathbf{h}_j) / \tau\bigr)
      + \sum_{k=1}^{K} \exp\!\bigl(-d_{\mathbb{H}}(\mathbf{h}_i, \mathbf{h}_{n_k}) / \tau\bigr)
    },
\end{equation}
where $\mathcal{E}_s$ is a sampled subset of edges, $\{n_k\}_{k=1}^K$ are negative samples, and $\tau$ is the temperature.

\paragraph{Cross-level contrastive learning.}
For each locality $\ell$, we pair it with a randomly chosen member segment $i \in \ell$ (positive) and $K$ random segments (negatives):
\begin{equation}\label{eq:loss_cross_cl}
  \mathcal{L}_{\text{cross}}
  = -\frac{1}{|\mathcal{L}_c|}
    \sum_{(\ell, i) \in \mathcal{L}_c}
    \log \frac{
      \exp\!\bigl(-d_{\mathbb{H}}(\mathbf{h}_\ell^{\text{loc}}, \mathbf{h}_i) / \tau\bigr)
    }{
      \exp\!\bigl(-d_{\mathbb{H}}(\mathbf{h}_\ell^{\text{loc}}, \mathbf{h}_i) / \tau\bigr)
      + \sum_{k=1}^{K} \exp\!\bigl(-d_{\mathbb{H}}(\mathbf{h}_\ell^{\text{loc}}, \mathbf{h}_{n_k}) / \tau\bigr)
    }.
\end{equation}

The combined contrastive loss is $\mathcal{L}_{\text{con}} = \mathcal{L}_{\text{seg}} + \mathcal{L}_{\text{cross}}$.

\subsubsection{Total Training Objective}\label{subsubsec:total_loss}

The overall loss is a weighted sum of all three components:
\begin{equation}\label{eq:total_loss}
  \mathcal{L}
  = \mathcal{L}_{\text{cls}}
  + \lambda_1\,\mathcal{L}_{\text{ent}}
  + \lambda_2\,\mathcal{L}_{\text{con}},
\end{equation}
where $\lambda_1$ and $\lambda_2$ balance the contributions of the entailment and contrastive objectives, respectively.
The entire model is trained end-to-end with the Adam optimiser and gradient clipping to prevent numerical instability inherent in hyperbolic operations.
